% Generado por sistema/backend ReporteTex.java -- no editar a mano.
% Datos: Base de datos: 148 importado (titulados) / Base de datos: 53 importado (empleadores)
\textit{Paso 1. Contar con la muestra y determinar su tamaño.} Respondieron la pregunta de situación laboral $n = 148$ titulados.

\textit{Paso 2. Contar la frecuencia de cada respuesta.}

\begin{mitabla}{Situación laboral actual ($n=148$)}
\begin{tabular}{lccc}
\toprule
Situación laboral & $f_i$ & $f_i / n$ & \% \\
\midrule
Trabaja en una organización & 108 & $108/148 = 0.730$ & 73.0 \\
Emprendimiento propio & 15 & $15/148 = 0.101$ & 10.1 \\
No trabaja & 25 & $25/148 = 0.169$ & 16.9 \\
\midrule
Total & 148 & & 100.0 \\
\bottomrule
\end{tabular}
\end{mitabla}

\textit{Lectura:} la respuesta más frecuente es ``Trabaja en una organización'' ($f_i = 108$, 73.0\% de 148).

\textit{Paso 3. Calcular la tasa de empleo.} Se consideran empleados quienes trabajan en una organización y quienes tienen un emprendimiento propio:
\[
\text{Tasa de empleo} = \frac{108 + 15}{148} = \frac{123}{148} = 0.8311 = 83.1\%
\]

